\documentclass[11pt]{article}

\usepackage[margin=1in]{geometry}
\usepackage{amsmath,amssymb,amsthm}
\usepackage{enumitem}
\usepackage[colorlinks=true,citecolor=blue,urlcolor=blue,linkcolor=blue]{hyperref}

\setlength{\parindent}{1.5em}
\setlength{\parskip}{0pt}

\newcommand{\R}{\mathbb R}
\newcommand{\E}{\mathbb E}
\newcommand{\Pp}{\mathbb P}
\newcommand{\G}{\Gamma}
\newcommand{\C}{\mathcal C}
\newcommand{\F}{\mathcal F}
\newcommand{\V}{\mathcal V}
\newcommand{\A}{\mathcal A}
\newcommand{\one}{\mathbf 1}
\newcommand{\ip}[2]{\langle #1,#2\rangle}
\newcommand{\norm}[1]{\left\|#1\right\|}
\newcommand{\Aut}{\operatorname{Aut}}
\newcommand{\Spec}{\operatorname{Spec}}
\newcommand{\Specsing}{\operatorname{Spec}_{\mathrm{sing}}}
\newcommand{\rank}{\operatorname{rank}}
\newcommand{\diam}{\operatorname{diam}}
\newcommand{\target}{\operatorname{target}}
\newcommand{\source}{\operatorname{source}}

\newtheorem{theorem}{Theorem}
\newtheorem{proposition}[theorem]{Proposition}
\newtheorem{lemma}[theorem]{Lemma}
\newtheorem{corollary}[theorem]{Corollary}
\newtheorem{definition}[theorem]{Definition}
\newtheorem{question}[theorem]{Question}
\newtheorem{conjecture}[theorem]{Conjecture}
\newtheorem{remark}[theorem]{Remark}

\title{Whitehead Field Dynamics and Reconfiguration Graphs of Cubic Multigraphs}
\author{Alexander Kolpakov}
\date{}

\begin{document}
\maketitle

\begin{abstract}
Let $\G_N$ be the graph whose vertices are isomorphism classes of connected
cubic multigraphs on $N$ vertices, with edges generated by Whitehead moves.  We
put over each class $\gamma=[G]$ the fibre $\F_\gamma=C^0(G)$ and transport a
scalar field by averaging the two endpoint values of the edge on which the
Whitehead move is performed.  This produces a random product of local
orthogonal projections, followed by settlement permutations.  The resulting
dynamics has a clean Lyapunov function, the variance above the constant mode,
and a noninvertible, semigroup-valued cocycle on a non-canonically trivialized
field bundle over $\G_N$.  We identify the exact configuration-model stationary
law of the local-move chain, prove a settlement-independent quenched exponent
upper bound of order $-N^{-3}$, and combine a spectral-average inequality with the
configuration-model Friedman theorem to obtain a stationary $-c/N$ upper
bound on the top norm Lyapunov exponent.  On the persistent labelled lift we
also find exact finite-time consensus, and hence exponent $-\infty$, at
$N=2,4,8$.  For the three graphs of Whitehead moves we derive a quadratic
family of low labelled modes, an exact loop-count Dirichlet form and sharp
bridge asymptotics for the multiplicity chain, and exact feature and field
obstructions for the simple support.  Gauge-covariant square defects remain
useful local observables, but immediate backtracking already shows why these
transports are not an ordinary connection.  Complete unlabelled enumeration
through $N=14$ separates proved asymptotics from finite conjectures.
\end{abstract}

\section{Introduction and contributions}

Fix an even integer $N$. Let $\C_N$ be the finite set of connected cubic
multigraphs on $N$ vertices, considered up to isomorphism unless labels are
explicitly introduced. Loops and multiple edges are allowed; a loop contributes
two to the degree. In addition to the carrier graph itself, there are three
graph-of-graphs conventions on this material. Li's labelled simple graph has
fixed-labelled cubic graphs as vertices and collapses repeated moves to the
same labelled target.  Grave de Peralta--Kolpakov's directed multigraph has
unlabelled classes $[G]$ as vertices and retains one directed slot for each
unlabelled Whitehead resolution, including multiplicity and self-slots.
Finally, the unlabelled simple-support graph has the same class vertices but
collapses all positive nonself multiplicities to one undirected edge. We use
$\G_N$ only when the accompanying text specifies the edge convention; every
degree, Laplacian, stationary law, and ball is indexed by its underlying
object.

The four state-and-clock conventions used below are summarized as follows.
{\small
\[
\begin{array}{c|c|c|c}
\text{object}&\text{states}&\text{one-step data}&\text{natural law}\\ \hline
\text{carrier field}&(G,f)&\text{one carrier edge and resolution}
    &\text{Rule A or specified Rule B}\\
\text{labelled simple GoG}&G\text{ on }[N]&\text{distinct labelled targets}
    &\text{degree-proportional}\\
\text{GDPK multiplicity chain}&[G]&3N\text{ resolution slots}
    &\pi_A\propto c(G)/|\Aut_v(G)|\\
\text{unlabelled simple support}&[G]&\text{distinct nonself targets}
    &\pi_B\propto q(G).
\end{array}
\]
}
The number of vertices of this GoG is not an unknown experimental count.  If
$N=2n$, then
\[
    |V(\G_N)|=a(n),
\]
where $a(n)$ is OEIS A005967, the number of isomorphism classes of connected
3-regular multigraphs of order $2n$, loops allowed \cite{oeis-a005967}.  The
first values are
\[
    2,5,17,71,388,2592,21096,\ldots
\]
for $N=2,4,6,8,10,12,14,\ldots$.  Thus enumeration code has a hard external
sanity check before any field dynamics is studied.

The papers of Grave de Peralta--Kolpakov and Li establish substantial
large-scale structure.  Grave de Peralta--Kolpakov prove that the
outer-conductance of $\G_N$ tends to zero, using the bottleneck given by
bridged cubic graphs \cite{gdp-kolpakov}.  Li shows that the graph of Whitehead
moves has increasingly many bounded Laplacian eigenvalues, using many local
cycle bottlenecks \cite{li}.  These results concern ordinary functions on the
vertex set of $\G_N$.

Here the moving object is different.  Over a cubic graph $G=(V,E)$ we place a
scalar field
\[
    f:V\to\R.
\]
So the fibre over the class $\gamma=[G]$ is
\[
    \F_\gamma=C^0(G)\simeq \R^{V(G)}.
\]
When a Whitehead move is performed on an edge $e=\{i,j\}$, the graph changes and
the field is updated by
\[
    f_i,f_j
    \mapsto
    \frac{f_i+f_j}{2},\frac{f_i+f_j}{2}.
\]
Afterwards one may resettle the new graph onto a chosen representative of its
isomorphism class; this only permutes field coordinates.

The field is scalar because no differential data is attached to the vertices
of $G$; the vector space $\R^N$ is only its coordinate representation.

\paragraph{Contributions.}
The paper makes four groups of distinctions and deductions.
\begin{enumerate}[leftmargin=*]
    \item For the carrier field we prove the exact heat identity, eliminate
    settlements by a pulled-back-pair factorization, obtain a uniform
    connectivity-only exponent bound, and prove exact finite-time consensus
    at $N=2,4,8$ on the persistent labelled lift.  We also isolate the
    carrier--field correlation that separates the true cocycle from an
    independently refreshed configuration benchmark.
    \item For Li's fixed-labelled simple graph we separate the simple,
    uniform-clock, and multiplicity walks.  We obtain asymptotic
    $L^1$-regularity, a quadratic-dimensional space of bounded
    combinatorial modes, fixed-radius ball asymptotics, and linear diameter
    and typical-distance lower bounds.
    \item For the Grave de Peralta--Kolpakov multiplicity chain we identify
    the exact configuration-model stationary law, spectral interval, Perron
    root, loop-count Dirichlet form, sharp bridge asymptotics, and explicit
    mixing and stationary-episode lower bounds.
    \item For the unlabelled simple support we give exact feature Rayleigh
    quotients and collision identities, prove conditional implications of
    degree growth, and characterize the instantaneous Rule-B field kernel by
    the active-edge carrier.
\end{enumerate}
The theoretical statements are separated below from complete finite
enumeration, conjectures, and the limited Lean verification layer.

\section{Whitehead Moves and the Exact Random Rule}

Let $e=\{u,v\}$ be a non-loop edge of a cubic graph.  At $u$ there are exactly
two other incident half-edges, call them $a,b$, and at $v$ there are exactly two
other incident half-edges, call them $c,d$.  Locally the edge $e$ sees the
partition
\[
    \{a,b\}\mid\{c,d\}.
\]
A Whitehead move collapses $e$ to a four-valent vertex and re-expands in one of
the two other pairings:
\[
    \{a,c\}\mid\{b,d\},
    \qquad
    \{a,d\}\mid\{b,c\}.
\]
On a fixed labelled vertex set, each of these two unordered pairings has two
assignments to the labelled central vertices $u,v$.  Thus there are four
labelled moves per non-loop edge, as in Li's labelled graph, but only two
unlabelled resolution slots per edge, as in Grave de Peralta--Kolpakov.  For a
cubic graph on $N$ vertices the corresponding raw counts are $6N$ and $3N$.
If $e$ is a loop, we use the convention of \cite{gdp-kolpakov}: the Whitehead
move is trivial.  In the half-edge model this convention is natural because the
pairing of half-edges is fixed and only the assignment of nearby half-edges to
vertices changes; when the two endpoints coincide there is no reassignment.

For the averaged field the factor of two can be compressed.  The two labelled
orientations of one unlabelled resolution differ by the transposition
$\tau=(uv)$, and
\[
    P_\tau A_{uv}=A_{uv}.
\]
With compatible settlements they therefore produce the same transported
field.  At a target with automorphisms they agree modulo a target
automorphism.  The transition table may consequently store two resolution
records per edge without losing the base-class kernel or the one-step
variance/source-Laplacian calculation.  For a full trajectory or a curvature
calculation, however, an intermediate target automorphism can change later
averages.  The compression is lossless there only after specifying coherent
settlements (or explicitly quotienting the fibre by stabilizers); otherwise
the omitted stabilizer holonomy is dynamical data.  The compression also does
not identify the two targets in the labelled simple graph and must not be used
for labelled ball counts.

There are two inequivalent random rules.

\begin{definition}[Rule A: local move first]
Start from a settled representative of the unlabelled class $\gamma=[G]$.
Choose a labelled Whitehead move upstairs, or equivalently for this averaged
field choose one of the two compressed unlabelled resolutions of an edge
object.  Apply this move to get a labelled graph $H$, project to the unlabelled
class $\delta=[H]$, and choose an isomorphism settling $H$ onto the chosen
representative of $\delta$.
\end{definition}

Rule A weights unlabelled transitions by concrete multiplicity.  If many
different edges and resolutions lead from $\gamma$ to the same class $\delta$,
then that edge of the unlabelled GoG is traversed more often.  This is the
natural rule if the physics is local in the current cubic graph.

\begin{definition}[Rule B: unlabelled neighbour first]
Start at $\gamma\in\G_N$, choose a neighbouring unlabelled class $\delta$
according to a rule on the quotient graph, and only afterwards choose a concrete
lift of the transition $\gamma\to\delta$.
\end{definition}

Rule B is a quotient-walk rule.  It is not equivalent to Rule A unless the
number of concrete lifts is constant across all unlabelled neighbours.  Under
Rule B, the edge of $G$ being averaged need not be uniformly distributed among
edge objects of $G$.  The expected field step is then governed by a weighted
Laplacian rather than the ordinary Laplacian of $G$.

This distinction controls both the base transition law and the field
operator.  Rule A
is local first: every formal edge/resolution slot in the current cubic graph is
equally likely, so every edge object is equally likely to be averaged.  Rule B
is quotient first: if ten concrete moves land in one quotient neighbour and one
concrete move lands in another, a uniform-neighbour Rule B gives those two
neighbours equal mass and then redistributes that mass over their lifts.  Thus a
single exceptional lift can receive much more probability under Rule B than it
does under Rule A.

For a fixed class $\gamma=[G]$, the Rule B field step is therefore recorded as
a weighted source Laplacian
\[
    L_\gamma^B=\sum_{e=\{i,j\}}p_\gamma^B(e)(e_i-e_j)(e_i-e_j)^T,
\]
where $p_\gamma^B(e)$ is the total Rule B probability that the chosen lift
averages edge $e$.  Rule A corresponds to the special case
\[
    p_\gamma^A(e)=\frac{1}{|E(G)|}=\frac{2}{3N}
\]
for every edge object.  Numerically, the right question is not simply whether
the base chains on $\G_N$ differ, but whether $L_\gamma^B$ differs
substantially from $(2/3N)L_G$ on the states that the chain actually visits.

\begin{remark}[What labels are doing]
The base graph $\G_N$ can stay unlabelled.  Labels on the vertices of the
current cubic graph are accounting coordinates for the field.  A step consists
of choosing a labelled lift, applying a concrete move, and then resettling onto
some representative of the target isomorphism class.  When automorphisms are
trivial, the settlement ambiguity essentially disappears once representatives
are fixed.  When automorphisms are nontrivial, the ambiguity is an orthogonal
permutation of the fibre, hence a gauge ambiguity rather than an extra
contraction mechanism.
\end{remark}

\begin{proposition}[Configuration-model stationary law of Rule A]
For a labelled adjacency multigraph $G$, let $m_{ii}$ be the number of loops at
$i$, let $m_{ij}$ be the number of edges between $i<j$, and put
\[
    c(G)=
    \frac{(3!)^N}
    {\prod_i2^{m_{ii}}m_{ii}!\prod_{i<j}m_{ij}!}.
\]
Let $\Aut_v(G)$ denote the vertex-permutation automorphism group.  If
$m_{GH}$ is the number of the $3N$ compressed Grave
de Peralta--Kolpakov slots from a representative of $[G]$ to the class $[H]$,
then
\[
    \frac{c(G)}{|\Aut_v(G)|}m_{GH}
    =
    \frac{c(H)}{|\Aut_v(H)|}m_{HG}.
\]
Consequently the Rule-A class chain
\[
    P_A([G],[H])=\frac{m_{GH}}{3N}
\]
is reversible with stationary law
\[
    \pi_A([G])
    =
    \frac{c(G)/|\Aut_v(G)|}
    {\sum_{[K]\in\G_N}c(K)/|\Aut_v(K)|}.
\]
This is exactly the cubic configuration-model law, conditioned on
connectedness and projected to unlabelled classes.
\end{proposition}

\begin{proof}
Give each labelled vertex three distinguishable half-edges.  The number of
perfect matchings of these half-edges which realize a fixed labelled adjacency
multigraph $G$ is $c(G)$: the denominator removes the two orientations of each
loop, permutations of loops at a vertex, and permutations of parallel edges.
An unlabelled class has $N!/|\Aut_v(G)|$ labelled adjacency realizations.

On fully labelled half-edge configurations, retain the four endpoint-oriented
moves at a non-loop edge (and four duplicate holding slots at a loop).  Each
move has a unique reverse, so uniform measure on configurations is reversible.
The four oriented moves collapse in pairs to the two unlabelled resolutions
recorded by $m_{GH}$.  Lumping the reversible configuration chain by
isomorphism class gives the displayed reciprocity identity and stationary law.
Whitehead moves preserve connectedness, so restricting and conditioning on
the connected component is legitimate.
\end{proof}

\begin{remark}
Thus stationarity is not a free assumption, and there is no third unknown
Rule-A graph distribution: its carrier marginal is the configuration model.
The weights $c(G)$ are essential when loops and parallel edges occur.  Already
for $N=2$, the triple-edge class and the class with one joining edge and one
loop at each endpoint have equal vertex-automorphism orders but stationary
masses $2/5$ and $3/5$.
\end{remark}

\section{Projection Calculus}

For a labelled vertex set $V=\{1,\ldots,N\}$ and a non-loop edge
$e=\{i,j\}$, define the local averaging operator
\[
    (A_{ij}f)_i=(A_{ij}f)_j=\frac{f_i+f_j}{2},
    \qquad
    (A_{ij}f)_k=f_k\quad(k\neq i,j).
\]
Let $e_i$ be the $i$th coordinate vector and put
\[
    u_{ij}=\frac{e_i-e_j}{\sqrt 2}.
\]
Then
\[
    A_{ij}=I-u_{ij}u_{ij}^T
    =
    I-\frac12(e_i-e_j)(e_i-e_j)^T.
\]

\begin{lemma}[Both pieces are projections]
$A_{ij}$ is the orthogonal projection onto the hyperplane
\[
    H_{ij}=\{f:f_i=f_j\}.
\]
The complementary operator
\[
    I-A_{ij}=\frac12(e_i-e_j)(e_i-e_j)^T
\]
is the orthogonal projection onto the line $\R(e_i-e_j)$.
\end{lemma}

\begin{proof}
The vector $u_{ij}$ is a unit normal to $H_{ij}$.  Therefore
$u_{ij}u_{ij}^T$ is the orthogonal projection onto $H_{ij}^\perp$ and
$I-u_{ij}u_{ij}^T$ is the orthogonal projection onto $H_{ij}$.  This also gives
$A_{ij}^2=A_{ij}$ and $(I-A_{ij})^2=I-A_{ij}$.
\end{proof}

Thus neither $A_{ij}$ nor $I-A_{ij}$ is nilpotent.  The contrast removed by the
averaging step is
\[
    (I-A_{ij})f
    =
    \frac{f_i-f_j}{2}(e_i-e_j).
\]
This expression does not require an orientation of the edge: exchanging $i$ and
$j$ changes both factors by a sign.

The basic noncommutativity is already visible on three vertices.  On
$\operatorname{span}\{e_1,e_2,e_3\}$,
\[
    [A_{23},A_{12}]
    =
    A_{23}A_{12}-A_{12}A_{23}
    =
    \frac14
    \begin{pmatrix}
        0&1&-1\\
        -1&0&1\\
        1&-1&0
    \end{pmatrix}.
\]
Consequently
\[
    \Specsing([A_{23},A_{12}])
    =
    \left\{\frac{\sqrt3}{4},\frac{\sqrt3}{4},0\right\}.
\]
This is the local algebraic seed of curvature: disjoint averages commute,
incident averages generally do not.

\begin{lemma}[Commutation criterion for elementary averages]
The projections $A_{ij}$ and $A_{k\ell}$ commute if the unordered pairs
$\{i,j\}$ and $\{k,\ell\}$ are equal or disjoint.  If they share exactly one
vertex, they do not commute.
\end{lemma}

\begin{proof}
If the pairs are equal this is trivial.  If they are disjoint, the operators
act on disjoint coordinate blocks.  For incident pairs it suffices, after
relabeling, to compute the commutator above, which is nonzero.
\end{proof}

\section{Variance and Heat Step}

Define the mean and variance by
\[
    \bar f=\frac1N\sum_{k=1}^N f_k,
    \qquad
    \V(f)=\sum_{k=1}^N(f_k-\bar f)^2
    =
    \norm{f-\bar f\one}^2.
\]
The constants are the consensus subspace.  The variance is the squared distance
from this subspace.

\begin{lemma}[One-step Lyapunov identity]
For one averaging step on $\{i,j\}$,
\[
    \sum_{k=1}^N(A_{ij}f)_k=\sum_{k=1}^Nf_k
\]
and
\[
    \V(A_{ij}f)=\V(f)-\frac12(f_i-f_j)^2.
\]
\end{lemma}

\begin{proof}
The sum of the $i,j$ coordinates is preserved.  Let $g=f-\bar f\one$.  Since
$A_{ij}\one=\one$, the mean-zero part after averaging is $A_{ij}g$.  Since
$A_{ij}$ is an orthogonal projection,
\[
    \norm{A_{ij}g}^2
    =
    \norm{g}^2-\ip{g}{u_{ij}}^2
    =
    \V(f)-\frac12(f_i-f_j)^2.
\]
\end{proof}

Let $L_G$ be the combinatorial Laplacian of the current cubic multigraph:
\[
    L_G=\sum_{\{i,j\}\in E(G),\,i\neq j}
    (e_i-e_j)(e_i-e_j)^T.
\]
Loops contribute zero to this sum.  Parallel edges are counted with
multiplicity.  If an edge object is chosen uniformly from all $|E(G)|=3N/2$
edge objects, then
\[
    \E[A_{ij}f\mid G,f]
    =
    \left(I-\frac{1}{2|E(G)|}L_G\right)f
    =
    \left(I-\frac{1}{3N}L_G\right)f,
\]
and
\[
    \E[\V(f')\mid G,f]
    =
    \V(f)-\frac{1}{2|E(G)|}f^TL_Gf
    =
    \V(f)-\frac{1}{3N}f^TL_Gf.
\]
If $\lambda_2(L_G)\ge c_0$, then Rayleigh--Ritz gives
\[
    \E[\V(f')\mid G,f]
    \le
    \left(1-\frac{c_0}{3N}\right)\V(f).
\]

For a non-uniform edge law $\mu$ on edge objects of $G$, the same calculation
uses the weighted Laplacian
\[
    L_{G,\mu}
    =
    \sum_{\{i,j\}\in E(G),\,i\neq j}
    \mu(\{i,j\})(e_i-e_j)(e_i-e_j)^T
\]
and gives
\[
    \E[\V(f')\mid G,f]
    =
    \V(f)-\frac12 f^TL_{G,\mu}f.
\]
Thus Rule B replaces the ordinary carrier Laplacian by a source-dependent
weighted Laplacian.

\section{The Field Bundle and Gauge}

The unlabelled graph-of-graphs has no canonical coordinates in its fibres.  The
right object is the bundle
\[
    \gamma=[G]\in\G_N,
    \qquad
    \F_\gamma=C^0(G).
\]
Choosing a representative $C(\gamma)$ for every class and choosing settlement
isomorphisms after every move is a choice of gauge.

For a concrete Whitehead move $m:\gamma\to\delta$, a settlement gives a linear
map
\[
    T_m:\F_\gamma\to\F_\delta,
    \qquad
    T_m=P_\sigma A_{ij},
\]
where $P_\sigma$ is the permutation matrix induced by the settling isomorphism.
If the source and target frames are changed by permutation matrices
$Q_\gamma,Q_\delta$, then the same transport is represented by
\[
    T_m\mapsto Q_\delta T_m Q_\gamma^{-1}.
\]
Thus matrix entries of $T_m$ are not intrinsic.  Norms, ranks, and singular
values of endpoint-covariant expressions are intrinsic.

\begin{lemma}[Gauge does not affect the Lyapunov identity]
For one settled unlabelled step
\[
    f'=P_\sigma A_{ij}f,
\]
the mean is preserved and
\[
    \V(f')=\V(f)-\frac12(f_i-f_j)^2.
\]
\end{lemma}

\begin{proof}
The permutation $P_\sigma$ preserves the constant vector, coordinate sums, and
Euclidean norm.  Apply the labelled variance identity before the permutation.
\end{proof}

This is the reason the one-step Lyapunov identity is clean: settlement choices
are orthogonal gauge changes, not hidden dissipation.  They are nevertheless
pathwise dynamical data, because cumulative settlements conjugate the endpoint
pairs averaged at later times.  The exact pullback formula below separates
these two statements.

\begin{proposition}[Exact settlement pullback]
Let $T_s=P_sA_{e_s}$, where $P_s$ is an arbitrary settlement permutation.  Put
\[
    Q_0=I,\qquad Q_{s+1}=P_sQ_s,
    \qquad \widetilde e_s=Q_s^{-1}e_s.
\]
Then
\[
    T_{t-1}\cdots T_0
    =Q_tA_{\widetilde e_{t-1}}\cdots A_{\widetilde e_0}.
\]
Consequently the singular values of a settled product are exactly those of a
product of ordinary pair averages in one persistent frame.
\end{proposition}

\begin{proof}
If $T_{s-1}\cdots T_0=Q_sY_s$, then
\[
    Y_{s+1}=Q_{s+1}^{-1}P_sA_{e_s}Q_sY_s
           =Q_s^{-1}A_{e_s}Q_sY_s
           =A_{Q_s^{-1}e_s}Y_s.
\]
Induction proves the factorization.
\end{proof}

\section{Covariant Difference and Square Curvature}

Let $s$ be a section of the field bundle over some portion of $\G_N$, so
$s_\gamma\in\F_\gamma$.  Given a concrete move $m:\gamma\to\delta$, define the
covariant difference
\[
    (Ds)_m=s_\delta-T_m s_\gamma.
\]
These maps are not invertible: every nonloop $T_m$ has rank $N-1$.  They
therefore form a dissipative semigroup cocycle rather than parallel transport
for a vector-bundle connection in the standard sense.  Comparing two paths
with the same endpoints is still meaningful, and squares give a useful local
comparison, but they are not the first obstruction to path independence.

Consider two concrete two-step paths with common source and target:
\[
    p:\gamma\to\alpha\to\delta,
    \qquad
    q:\gamma\to\beta\to\delta.
\]
After choosing local settlements, put
\[
    T_p=T_{p_2}T_{p_1},
    \qquad
    T_q=T_{q_2}T_{q_1},
\]
and define the square defect
\[
    \Omega(p,q)=T_p-T_q.
\]

\begin{proposition}[Gauge covariance of the square defect]
Under a coherent gauge change at the common source and common target,
\[
    \Omega(p,q)\mapsto Q_\delta\Omega(p,q)Q_\gamma^{-1}.
\]
Consequently
\[
    \Specsing(T_p-T_q),
    \quad
    \norm{T_p-T_q}_{\mathrm{op}},
    \quad
    \norm{T_p-T_q}_{\mathrm{Frob}},
    \quad
    \rank(T_p-T_q)
\]
are independent of the chosen source and target frames.
\end{proposition}

\begin{proof}
Both path transports transform by the same endpoint permutations:
\[
    T_p\mapsto Q_\delta T_p Q_\gamma^{-1},
    \qquad
    T_q\mapsto Q_\delta T_q Q_\gamma^{-1}.
\]
Subtract.  Left and right multiplication by orthogonal permutation matrices
preserves singular values, operator norm, Frobenius norm, and rank.
\end{proof}

Thus the clean curvature definition is not a matrix of entries.  It is the
singular spectrum of the square defect:
\[
    \operatorname{Curv}(p,q)
    =
    \Specsing(T_p-T_q),
\]
or, as a scalar,
\[
    \kappa(p,q)=\norm{T_p-T_q}_{\mathrm{op}}.
\]
The elementary operators $A_{ij}$ and $I-A_{ij}$ are projections, not
differentials.  The displayed $D$ is currently defined only from vertices to
edges.  Calling the square defect the obstruction to $D^2=0$ would require a
separate extension of $D$ from edges to chosen $2$-cells; without that extra
definition, the justified statement is gauge covariance of the path
difference.

There is already a one-dimensional backtracking defect.  If a Whitehead move
on $e$ is followed by its graph inverse with coherent inverse settlement, then
in the original frame
\[
    T_{m^{-1}}T_m=A_e^2=A_e\ne I,
\]
and $I-A_e$ is a rank-one orthogonal projection of operator norm one.  Thus
even a contractible move--inverse path is dissipative.

\begin{proposition}[Flatness for disjoint moves]
If two concrete Whitehead moves are performed on non-loop edges whose endpoint
sets are disjoint, then the corresponding labelled square defect is zero.  Any
coherently settled unlabelled square has
\[
    \Specsing(T_p-T_q)=\{0,\ldots,0\}.
\]
\end{proposition}

\begin{proof}
In the half-edge model, each Whitehead move only reassigns the four half-edges
incident to the endpoints of the selected edge and not belonging to that edge.
Disjoint endpoint sets give disjoint local data, so the graph rewrites commute.
The field operators also commute because the averaging projections act on
disjoint coordinate pairs.  Thus both two-step path transports are identical,
up to the same final settlement.
\end{proof}

Incident squares are the basic non-flat squares.  Their defect is not only a
property of the cubic graph rewrite; it also includes the averaging operators.
Locally, when the two consecutive averaged pairs share one vertex, the
commutator has singular values
\[
    \frac{\sqrt3}{4},\frac{\sqrt3}{4},0
\]
on the three affected coordinates.  In an actual GoG square this local algebra
is dressed by the two graph rewrites and the endpoint settlements.

Automorphisms create one extra subtlety.  If the endpoint graph has a nontrivial
automorphism group, two apparently different final settlements may differ by
$U\in\Aut(\delta)$.  A quotient-level aligned defect is therefore
\[
    \kappa_{\mathrm{orb}}(p,q)
    =
    \min_{U\in\Aut(\delta)}
    \norm{T_p-UT_q}_{\mathrm{op}}.
\]
This is not needed when the two lifted paths and settlements are tracked
coherently, but it is useful when comparing defects computed from canonical
representatives.  Averaging over all settlements,
\[
    \overline T_m
    =
    \frac{1}{|\Aut(\delta)|}
    \sum_{\alpha\in\Aut(\delta)}P_\alpha T_m,
\]
is canonical but changes the dynamics: it projects toward
automorphism-invariant components and can introduce artificial contraction.

\section{Deterministic Consensus from Products of Projections}

Before considering random graph motion, isolate the deterministic fact about
averaging projections.  Let
\[
    E_0=\one^\perp=\{f\in\R^N:\sum_i f_i=0\}.
\]
Every $A_{ij}$ preserves $E_0$ and has norm at most one on $E_0$.

\begin{lemma}[Equality case]
Let $P=A_{i_rj_r}\cdots A_{i_1j_1}$.  If $f\in E_0$ and
\[
    \norm{Pf}=\norm{f},
\]
then $f_{i_s}=f_{j_s}$ for every averaged pair $\{i_s,j_s\}$.
\end{lemma}

\begin{proof}
Each $A_{i_sj_s}$ is an orthogonal projection, so norms can only decrease at
each step.  Equality at the end forces equality at every intermediate step.
For the step $A_{ij}$, equality in the projection identity is exactly
$f_i=f_j$ for the vector before that step.  Since previous projections do not
increase norm, the same conclusion propagates along the equality chain.
\end{proof}

\begin{corollary}[Exact connected-block criterion]
Let $H$ be the graph on $\{1,\ldots,N\}$ whose edges are the averaged pairs
$\{i_s,j_s\}$.  Then
\[
    \norm{A_{i_rj_r}\cdots A_{i_1j_1}|_{E_0}}<1
    \quad\Longleftrightarrow\quad H\text{ is connected}.
\]
\end{corollary}

\begin{proof}
If $H$ is disconnected, a nonzero mean-zero field constant on its components
is fixed by every factor.  Conversely, if equality held for some nonzero
$f\in E_0$, the lemma would force $f$ to be constant on every component of
$H$.  Connectedness would make it globally constant, hence zero, a
contradiction.
\end{proof}

\begin{proposition}[Quantitative connected-block contraction]
Let $m$ pair averages have connected union graph, and put
$a_N=2(1-\cos(\pi/N))$.  Then
\[
 \norm{A_{i_mj_m}\cdots A_{i_1j_1}|_{E_0}}^2
 \le
 1-\frac{a_N}{2(N-1)(1+\sqrt m)^2}.
\]
Consequently, if every consecutive block of $m$ pulled-back effective pairs
is connected, arbitrary settlements still give deterministic exponential
consensus, with exponent at most
\[
 \frac1{2m}\log\left(
 1-\frac{a_N}{2(N-1)(1+\sqrt m)^2}\right)<0.
\]
\end{proposition}

\begin{proof}
Let $x^s$ be the field after $s$ averages,
$d_s=x^{s-1}_{i_s}-x^{s-1}_{j_s}$, and
\[
 D=\V(x^0)-\V(x^m)=\frac12\sum_s d_s^2.
\]
The total displacement before any step is at most
\[
 \sum_s\norm{x^s-x^{s-1}}
 =\sum_s\frac{|d_s|}{\sqrt2}\le\sqrt{mD}.
\]
Choose one occurrence of each edge of a spanning tree of the union graph.
At the time such an edge $ij$ is averaged,
\[
 |x_i^0-x_j^0|
 \le |d_s|+\sqrt2\norm{x^{s-1}-x^0}
 \le\sqrt{2D}(1+\sqrt m).
\]
The spanning-tree Poincar\'e inequality and the path-minimality of algebraic
connectivity give
\[
 a_N\V(x^0)
 \le\sum_{ij\in T}(x_i^0-x_j^0)^2
 \le2(N-1)(1+\sqrt m)^2D.
\]
Rearranging proves the one-block estimate.  Multiplying the estimates over
disjoint blocks proves the exponent bound.
\end{proof}

Thus the relevant connectivity condition is imposed on the endpoint pairs
in one persistent frame, rather than on physical source edges written in
successively changing frames.

\section{\texorpdfstring{Finite-$N$}{Finite-N} Exponential Consensus}

The following theorem gives a robust fixed-$N$ baseline.  It says nothing sharp
about the constant as $N$ grows, and at some sizes the conclusion strengthens
to finite-time exact consensus.

\begin{theorem}[Fixed-$N$ exponential consensus on the labelled lift]
Fix $N$.  Let $(X_t)$ be a finite irreducible Markov chain on labelled connected
cubic multigraphs on the persistent vertex set $\{1,\ldots,N\}$, and assume
that every concrete Whitehead move available from a state has positive
transition probability.  Do not resettle between steps, and let
\[
    B_t=A_{i_tj_t}
\]
be the field operator at time $t$, restricted to $E_0$.  Then there
exist $M\ge1$, $q>0$, and $\rho<1$ such that almost surely
\[
    \limsup_{t\to\infty}
    \frac1t\log\norm{B_{t-1}\cdots B_0|_{E_0}}
    \le
    \frac{q}{M}\log\rho
    <0.
\]
Consequently every initial field converges exponentially to its mean.
\end{theorem}

\begin{proof}
Choose one labelled connected cubic graph $G^\ast$ and a spanning tree
$T^\ast\subset G^\ast$.  Since the labelled chain is finite and irreducible,
from any state $x$ there is a path of Whitehead moves leading to $G^\ast$.  By
finiteness, the length of such a path can be chosen at most $D$ uniformly in
$x$, after padding if necessary.  Once at $G^\ast$, for each edge of
$T^\ast$ perform a chosen Whitehead move and immediately perform its inverse.
This averages every tree edge and returns the labelled carrier to $G^\ast$
after each pair.  Let $M$ be a
common upper bound for the total length of this ``success block''.

Full support and finiteness imply that the conditional probability of realizing
one of these success blocks in the next $M$ steps is bounded below by some
$q>0$, uniformly over the current state.

For every success block, the product of its averaging projections includes all
edges of the spanning tree $T^\ast$.  By the connected-block corollary, its
restriction to $E_0$ has norm strictly less than one.  There are only finitely
many success-block words of length at most $M$, so the maximum of their norms is
some $\rho<1$.

Now divide time into disjoint blocks of length $M$.  Let $Y_k$ be the indicator
that block $k$ is successful and put
$p_k=\E[Y_k\mid\mathcal F_{kM}]\ge q$.  The martingale
$\sum_{k<r}(Y_k-p_k)$ has bounded increments, so its strong law
\cite{williams,durrett} gives
\[
    \liminf_{r\to\infty}\frac1r\sum_{k<r}Y_k\ge q
    \quad\text{almost surely}.
\]
Successful blocks contract by at most $\rho$; all other blocks have norm at
most one.  Hence
\[
    \limsup_{r\to\infty}
    \frac1{rM}\log\norm{B_{rM-1}\cdots B_0|_{E_0}}
    \le
    \frac{q}{M}\log\rho.
\]
The leftover times between multiples of $M$ contribute only norm-one factors,
so the same limsup holds for all $t$.
\end{proof}

\begin{remark}
The theorem is deliberately stated on the persistent labelled lift.  For an
unlabelled implementation with intermediate settlements
$P_{\sigma_t}A_{i_tj_t}$, orthogonality of each permutation alone does not
preserve the connected-block argument: cumulative settlements conjugate the
later averaging pairs.  The same proof works if the averaged pairs are pulled
back to one persistent frame and those effective pairs form a connected block.
Without that condition, fixed-$N$ consensus for an arbitrary settlement rule
has not been proved here.
\end{remark}

\begin{theorem}[Finite-time consensus at $N=2,4,8$]
Under the hypotheses of the persistent labelled theorem, the restriction of
the field product to $E_0$ becomes exactly zero in finite time almost surely at
$N=2,4,8$.  Thus its top exponent is $-\infty$ at these sizes.
\end{theorem}

\begin{proof}
A Whitehead move on a carrier edge $e$, immediately followed by its graph
inverse, returns the labelled carrier and applies $A_e^2=A_e$ to the field.
It therefore embeds any finite edge-average word supported by one carrier into
a returning Whitehead word.

At $N=2$, one nonloop average is consensus.  At $N=4$, take $K_4$ and average
the two perfect matchings
\[
    \{12,34\},\qquad \{13,24\}.
\]
The first round makes two equal-value pairs and the second merges them.  At
$N=8$, take the cube $Q_3$ and average its three coordinate perfect matchings;
the equal-value blocks double from sizes $1$ to $2,4,8$.  In each case the
product is exactly $J=N^{-1}\one\one^T$.

The finite irreducible carrier chain visits the chosen carrier infinitely
often.  Full support gives a fixed positive probability of following the
finite killing word at each visit, so the word occurs almost surely.  Every
later doubly stochastic factor fixes $J$.
\end{proof}

\begin{proposition}[Dyadic obstruction]
A finite product of pair averages and permutations can equal $J$ only if $N$
is a power of two.  If $N$ is not a power of two, then after $t$ averaging
steps
\[
    \norm{T_{t-1}\cdots T_0|_{E_0}}
    =\norm{T_{t-1}\cdots T_0-J}
    \ge \frac{1}{N2^t}.
\]
\end{proposition}

\begin{proof}
Every entry of the product is dyadic with denominator dividing $2^t$.  Killing
$E_0$ while fixing constants forces the full product to be $J$.  Its entries
$1/N$ are dyadic only when $N$ is a power of two.  Otherwise some entry differs
from $1/N$ by at least $1/(N2^t)$, and the operator norm dominates every
matrix entry.
\end{proof}

The proof also explains why the bridged bottleneck does not automatically
destroy finite-$N$ field relaxation.  A successful block may pass through
bridged graphs and may average bridge edges.  That is fine: the block contracts
because the list of averaged label-pairs contains a spanning tree, not because
the intermediate base graphs avoid bottlenecks.

\section{Pathwise Lyapunov Bounds from Spectral Expansion}

The fixed-$N$ theorem gives negativity but not the scale.  The scale expected
from a random local averaging step on an expander cubic graph is $-c/N$.

There is first a completely uniform bound which does not assume expansion or
stationarity.

\begin{theorem}[Connectivity-only quenched bound]
For every connected cubic multigraph $G$ on $N$ vertices,
\[
    \lambda_2(L_G)\ge
    a_N:=2\bigl(1-\cos(\pi/N)\bigr).
\]
Consequently every adapted trajectory of connected carriers under uniform-edge
Rule A, with arbitrary settlement permutations, satisfies
\[
    \limsup_{t\to\infty}\frac1t
    \log\norm{T_{t-1}\cdots T_0|_{E_0}}
    \le \frac12\log\left(1-\frac{a_N}{3N}\right)
    =-\frac{\pi^2}{6N^3}+O(N^{-5}).
\]
No irreducibility, stationary initial law, or expander frequency is required.
\end{theorem}

\begin{proof}
Delete loops and choose a spanning tree $T$ of $G$, retaining one edge from a
parallel family.  Then $L_G\succeq L_T$.  The path minimizes algebraic
connectivity among trees on $N$ vertices, and its gap is
$2(1-\cos(\pi/N))$.

The conditional heat identity now gives
\[
  \E[\V(f_{t+1})\mid\mathcal F_t]
  \le\left(1-\frac{a_N}{3N}\right)\V(f_t).
\]
Settlement preserves variance.  Dividing by the deterministic geometric
factor produces a nonnegative supermartingale.  Applying its almost-sure
bound to a basis of $E_0$ gives the operator statement.
\end{proof}

The worst-carrier order cannot be improved from connectivity alone.  Put a
loop at each endpoint of a path and give consecutive path edges alternating
multiplicities $1,2,1,2,\ldots,1$.  This is cubic, and the linear test field
has Rayleigh quotient
\[
    \frac{18N-24}{N(N^2-1)}=O(N^{-2}).
\]
Thus the uniform $N^{-3}$ exponent and the typical stationary $N^{-1}$
exponent answer different questions.

Let $(G_t,f_t)$ be a Rule A process in which, conditional on $G_t$, the averaged
edge object is uniform among the $3N/2$ edge objects of $G_t$.  Let
\[
    H_t=\{\lambda_2(L_{G_t})\ge c_0\},
    \qquad
    \alpha_N=\frac{c_0}{3N},
\]
where $0<c_0<3N$.

\begin{theorem}[Good-graph frequency implies a $-c/N$ field exponent]
Assume that for some $p\in[0,1]$,
\[
    \liminf_{t\to\infty}
    \frac1t\sum_{s=0}^{t-1}\mathbf 1_{H_s}
    \ge p
    \qquad\text{almost surely}.
\]
Then for every initial field $f_0$,
\[
    \limsup_{t\to\infty}
    \frac1t\log\V(f_t)
    \le
    p\log(1-\alpha_N)
    \qquad\text{almost surely}.
\]
Equivalently, if $f_0$ is not already constant,
\[
    \limsup_{t\to\infty}
    \frac1t\log\norm{f_t-\bar f_0\one}
    \le
    \frac p2\log(1-\alpha_N)
    =
    -\frac{pc_0}{6N}+O(N^{-2}).
\]
For fixed $N$, the same upper bound holds for the top operator exponent on
$E_0$.
\end{theorem}

\begin{proof}
The conditional heat-step identity and the Rayleigh bound imply
\[
    \E[\V(f_{t+1})\mid\mathcal F_t]
    \le
    (1-\alpha_N\mathbf 1_{H_t})\V(f_t).
\]
Define
\[
    Q_t=\prod_{s=0}^{t-1}(1-\alpha_N\mathbf 1_{H_s}),
    \qquad
    Z_t=\frac{\V(f_t)}{Q_t}.
\]
Since $H_t$ is determined by $G_t$, the factor $Q_{t+1}/Q_t$ is
$\mathcal F_t$-measurable.  Hence $(Z_t)$ is a nonnegative supermartingale and
therefore converges almost surely to a finite limit.  Thus
\[
    \limsup_{t\to\infty}\frac1t\log\V(f_t)
    \le
    \limsup_{t\to\infty}\frac1t\log Q_t.
\]
But
\[
    \frac1t\log Q_t
    =
    \frac1t\sum_{s=0}^{t-1}
    \mathbf 1_{H_s}\log(1-\alpha_N),
\]
and $\log(1-\alpha_N)<0$.  The assumed lower frequency of $H_s$ gives the stated
upper bound.

The norm bound is half the variance bound.  For the operator norm on $E_0$,
apply the vector bound to a fixed orthonormal basis of $E_0$ and intersect the
finitely many almost-sure events.  If every basis vector decays at rate at most
$r$, then every vector decays at rate at most $r$ up to a fixed dimension
constant, which disappears after division by $t$.
\end{proof}

\begin{theorem}[Spectral-average version]
Assume $N\ge4$.  For the same Rule A process, set
\[
    \ell_t=\lambda_2(L_{G_t}).
\]
If
\[
    \liminf_{t\to\infty}\frac1t\sum_{s=0}^{t-1}\ell_s\ge \ell_*
    \qquad\text{almost surely},
\]
then
\[
    \limsup_{t\to\infty}\frac1t\log\V(f_t)
    \le
    \limsup_{t\to\infty}\frac1t
    \sum_{s=0}^{t-1}\log\left(1-\frac{\ell_s}{3N}\right)
    \le
    -\frac{\ell_*}{3N},
\]
and therefore
\[
    \limsup_{t\to\infty}
    \frac1t\log\norm{f_t-\bar f_0\one}
    \le
    -\frac{\ell_*}{6N}.
\]
\end{theorem}

\begin{proof}
The conditional variance estimate holds with the actual current spectral gap:
\[
    \E[\V(f_{t+1})\mid\mathcal F_t]
    \le
    \left(1-\frac{\ell_t}{3N}\right)\V(f_t).
\]
The factor is positive because the Laplacian eigenvalues of a cubic multigraph
lie in $[0,6]$ and $3N>6$.  (At $N=2$, a zero factor means one-step consensus
and can be handled separately.)  Define
\[
    Q_t=\prod_{s=0}^{t-1}\left(1-\frac{\ell_s}{3N}\right).
\]
As before, $\V(f_t)/Q_t$ is a nonnegative supermartingale.  This gives the
first displayed bound.  The second follows from
\[
    \log(1-x)\le -x
    \qquad (0\le x<1).
\]
\end{proof}

This gives precise meaning to saying that the current graph is an expander on
average.  What is needed is not stationarity itself, but a pathwise lower
bound on the time average of the cubic-graph spectral gap seen by the moving
field.

\begin{corollary}[Stationary law as one source of good frequency]
Suppose the graph component is a finite irreducible Markov chain with
stationary law $\pi_N$.  If
\[
    \pi_N\{G:\lambda_2(L_G)\ge c_0\}\ge 1-\varepsilon_N,
\]
then the theorem applies with $p=1-\varepsilon_N$, for every initial graph.
\end{corollary}

\begin{proof}
For a finite irreducible Markov chain, the ergodic theorem says that for every
function $\varphi$ on the state space,
\[
    \frac1t\sum_{s=0}^{t-1}\varphi(G_s)
    \to
    \sum_G\pi_N(G)\varphi(G)
    \quad\text{almost surely}.
\]
Apply this to the indicator of the good set.
\end{proof}

For concrete Rule A, the preceding stationary law is not abstract: it is the
cubic configuration model conditioned on connectedness.  Spectral typicality
can therefore be imported from the matching model itself rather than posed as
a new property of an unknown stationary law.

This corollary is not the only route.  Stationarity is convenient, but the
Lyapunov theorem only needs a pathwise lower frequency of good graphs.  A
statement such as ``the graph is an expander on average'' is not by itself
enough unless it implies such a frequency or a comparable supermartingale
inequality.

The Jensen/log heuristic is consistent with the supermartingale proof.  Whenever
$\V(f_t)>0$,
\[
    \E\left[
        \log\frac{\V(f_{t+1})}{\V(f_t)}
        \,\middle|\,\mathcal F_t
    \right]
    \le
    \log
    \E\left[
        \frac{\V(f_{t+1})}{\V(f_t)}
        \,\middle|\,\mathcal F_t
    \right].
\]
On a good graph this is at most $\log(1-c_0/(3N))$.  The supermartingale
argument is stronger because it converts conditional expected contraction into
an almost-sure pathwise upper bound.

\section{Exact One-Step Identities Beyond Variance}

The configuration marginal gives a useful reset benchmark, but it must not be
mistaken for the correlated trajectory law.

\begin{proposition}[Independent configuration reset]
Let $G$ be an unconditioned cubic configuration, let $e$ be a uniform matched
edge object, and let a deterministic mean-zero $f$ be independent of $G$.
Then exactly
\[
    \E\V(A_ef)=\left(1-\frac{3}{3N-1}\right)\V(f),
\]
or equivalently
\[
    \E L_G=\frac{9N}{3N-1}
       \left(I-\frac1N\one\one^T\right).
\]
If $\ell_N=\E[\#\text{ loops}\mid G\text{ connected}]$, the connected reset
factor is
\[
    1-\frac{3N-2\ell_N}{3N(N-1)}.
\]
\end{proposition}

\begin{proof}
Expose an oriented uniform stub.  Conditional on its vertex being $i$, each
$j\ne i$ is the partner vertex with probability $3/(3N-1)$.  Hence
\[
 \E(f_i-f_j)^2
 =\frac{3}{N(3N-1)}\sum_i\sum_{j\ne i}(f_i-f_j)^2
 =\frac{6}{3N-1}\V(f),
\]
and one average loses half of this.  Under connected conditioning,
permutation invariance makes the mean Laplacian scalar on $E_0$; its trace
$3N-2\ell_N$ determines the scalar.
\end{proof}

Fresh independent resampling would iterate this factor and suggest norm
constant $1/2$ on the $N^{-1}$ scale.  The actual field records previously
selected edges and is correlated with $G_t$, so the identity does not iterate
along a Whitehead trajectory.

\begin{theorem}[Resolution-independent target energy]
Suppose a move on $e=\{i,j\}$ produces either raw resolution $G'$.  Then
\[
    (A_ef)^TL_{G'}(A_ef)=(A_ef)^TL_G(A_ef).
\]
For $b_e=e_i-e_j$ (and $b_e=0$ on a loop), put
\[
 \kappa_e=b_e^TL_Gb_e,\qquad
 D_G=\sum_e\kappa_e b_eb_e^T,\qquad M=\frac{3N}{2}.
\]
Uniform-edge Rule A therefore satisfies
\[
 \E[f'^TL_{G'}f'\mid G,f]
 =f^T\left(L_G-\frac1M L_G^2+\frac1{4M}D_G\right)f.
\]
For a simple cubic carrier, $D_G=8L_G$.
\end{theorem}

\begin{proof}
After $A_e$, the values at $i,j$ agree.  A Whitehead resolution only changes
which of these equal-valued endpoints owns each surrounding half-edge, so it
does not change any contrast energy.  Averaging $A_eL_GA_e$ over edge objects
and using $\sum_e b_eb_e^T=L_G$ gives the matrix identity.  In a simple cubic
graph, $b_e^TL_Gb_e=3+3+2=8$.
\end{proof}

Thus, for a simple carrier and an $L_G$-eigenfield of eigenvalue $\lambda$,
the expected carrier energy is multiplied by $1+(2-\lambda)/M$.  Low-frequency
energy can increase even while variance decreases; the next moment involves
$L_G^2$ and exposes the carrier--field alignment problem.

\begin{proposition}[Exact cut drift and cut count]
Let $|S|=k$, $b=|\partial_GS|$, and
$f=\one_S-(k/N)\one$.  Then
\[
 \frac{\E[\V(f')\mid G,f]}{\V(f)}
 =1-\frac{b}{3k(N-k)}.
\]
In an unconditioned cubic configuration, the expected number of $k$-vertex
sets with exactly $b$ crossing edge objects is
\[
 \binom Nk
 \frac{\binom{3k}{b}\binom{3(N-k)}{b}b!
 (3k-b-1)!!(3(N-k)-b-1)!!}{(3N-1)!!},
\]
when $3k-b$ is even, and zero otherwise.  For fixed $B$ and
$\varepsilon>0$, the probability of a cut satisfying
\[
 \varepsilon N\le |S|\le N/2,\qquad |\partial S|\le B,
\]
is at most $e^{-c_{\varepsilon,B}N}$ for all sufficiently large $N$; this
persists after conditioning on connectedness.
\end{proposition}

\begin{proof}
Only a crossing edge changes the two-level field, losing $1/2$ unit of
variance; divide the expected loss $b/(3N)$ by $k(N-k)/N$.  For the count,
choose the $b$ crossing stubs on both sides, match them bijectively, and pair
the remaining stubs internally.  Stirling's formula gives exponent
$-NH(k/N)/2+o(N)$ uniformly on the displayed balanced range.
\end{proof}

\section{Random Cubic Expansion and GoG Bottlenecks}

For a simple random cubic graph, Friedman's theorem gives, with probability
tending to one,
\[
    \lambda_2(L_G)
    \ge
    3-2\sqrt2-\eta
\]
for every fixed $\eta>0$ \cite{friedman}.  More directly for present purposes,
Bordenave proves the corresponding non-backtracking bound for a uniform
perfect matching of the cubic half-edges \cite{bordenave}.  The Ihara--Bass
relation for regular multigraphs gives
\[
    \lambda_2(L_G)\ge 3-2\sqrt2-o(1)
\]
with high probability in the cubic configuration model.  Random cubic
configurations are connected with high probability, so conditioning on
connectedness does not change this conclusion.  Thus the number
\[
    3-2\sqrt2\approx 0.1716
\]
is the natural benchmark lower gap for a typical Rule-A carrier.

\begin{corollary}[Unconditional Rule-A Lyapunov upper scale]
For any specified adapted settlement rule, define the upper top norm exponent
by
\[
 \Lambda_N^A:=\limsup_{t\to\infty}\frac1t
 \log\norm{T_{t-1}\cdots T_0|_{E_0}}.
\]
For the stationary Rule-A field process, and from every initial graph by
ergodicity,
\[
    \Lambda_N^A
    \le
    -\frac{3-2\sqrt2-o(1)}{6N}
\]
for the top norm exponent.
\end{corollary}

\begin{proof}
The configuration-model stationary identity and the preceding spectral theorem
give good-graph frequency $1-o(1)$.  Apply the good-graph-frequency theorem
with any $c_0<3-2\sqrt2$ and then let $c_0$ increase to the benchmark.
\end{proof}

\begin{proposition}[Sharp GDPK bridge flow for Rule A]
Let $B_N$ be the bridged classes, $B_0=B_N^c$, and $B_1$ the classes with
exactly one bridge.  For class sets $S,T$, write
\[
 Q_A(S,T)=\sum_{\substack{G\in S\\H\in T}}\pi_A(G)P_A(G,H),
\]
and normalize a cut by its smaller stationary side:
\[
 \Phi_A(B_N)=
 \frac{Q_A(B_N,B_0)}{\min\{\pi_A(B_N),\pi_A(B_0)\}}.
\]
Then exactly
\[
    Q_A(B_N,B_0)=\frac{2}{3N}\pi_A(B_1).
\]
Moreover,
\[
 \pi_A(B_0)\to e^{-1},\qquad
 \pi_A(B_1)\to e^{-1},\qquad
 \pi_A(B_N)\to1-e^{-1},
\]
and hence
\[
    \Phi_A(B_N)=\frac{2}{3N}(1+o(1)).
\]
In particular the Rule-A gap is $O(N^{-1})$ and the relaxation time is
$\Omega(N)$.
\end{proposition}

\begin{proof}
A move can cross from the bridged to the bridgeless side only by acting on the
unique bridge.  Conversely, let $e=uv$ be the unique bridge and let $A,C$ be
the two components of $G-e$.  Neither component has an internal bridge.  In
either resolution, one side edge at $u$ and one side edge at $v$ change
endpoint, and the two connected sides are then joined by those two edges and
$e$.  Each of these three cross edges lies on a cycle using the other two and
paths inside $A,C$; cycles internal to either side are unchanged or repaired
through $e$.  Thus both compressed resolutions are bridgeless.  A class in
$B_1$ has exactly two exiting slots, while a class with more than one bridge
has none.  Averaging under $\pi_A$ proves the exact flow formula.

It remains to determine the masses.  In the unconditioned configuration
model, the loop count has mean $3N/(3N-1)$ and converges by factorial moments
to $\operatorname{Pois}(1)$.  If a bridge cuts off a smaller vertex set of
size $k$, then $3k=2|E(S)|+1$, so $k$ is odd.  The expected number of
$k$-sets with exactly one crossing matched pair is
\[
 A_{N,k}=\binom Nk(3k)(3N-3k)
 \frac{(3k-2)!!(3N-3k-2)!!}{(3N-1)!!}.
\]
This is a first-moment upper bound because the two internally paired pieces
need not be connected.  The first term is
\[
 A_{N,3}=\frac{35}{18N}+O(N^{-2}).
\]
Exact double-factorial ratios give
$A_{N,k}=O(N^{-(k-1)/2})$ for each fixed odd $k$, and a uniform Stirling bound
makes the middle range exponentially small.  Consequently
\[
 \sum_{\substack{3\le k\le N/2\\k\ {\rm odd}}}A_{N,k}
 =\frac{35}{18N}+O(N^{-2})=O(N^{-1}).
\]
Conditioning on connectedness changes these estimates by $1+O(N^{-1})$.
Indeed, the analogous expected count of an isolated even $k$-set is
\[
 D_{N,k}=\binom Nk
 \frac{(3k-1)!!(3N-3k-1)!!}{(3N-1)!!};
\]
its $k=2$ term is $5/(18N)+O(N^{-2})$, and the same ratio and Stirling
estimates give total disconnection probability $O(N^{-1})$.

A loop vertex has one remaining incident edge, which is a singleton-side
bridge.  Up to an event of probability $O(N^{-1})$, all bridges are of this
form.  Hence $B_0$ differs from $\{L=0\}$ by probability $O(N^{-1})$ and
$B_1$ differs from $\{L=1\}$ by probability $O(N^{-1})$.  The Poisson limit
gives the three stationary-mass limits, and substitution in the exact flow
formula gives the conductance asymptotic.
\end{proof}

Starting from stationarity conditioned on $B_0$ and testing total variation on
$B_0$ gives, for fixed $\varepsilon<1-e^{-1}$,
\[
 t_{\rm mix}(\varepsilon)
 \ge \frac{3N}{2}(1-e^{-1}-\varepsilon)(1+o(1)).
\]
by the standard bottleneck argument \cite{levin-peres-wilmer}.
Thus $t_{\rm mix}(1/4)\ge0.57318\ldots N(1+o(1))$.  The stationary mean
bridgeless and bridged run lengths are respectively
\[
 \frac{3N}{2}(1+o(1)),\qquad
 \frac{3N}{2}(e-1)(1+o(1)).
\]
\begin{conjecture}[Finite bridge identity]
For every even $N\ge4$,
\[
 \Phi_A(B_N)=\frac1{3N/2-1}.
\]
Equivalently, if $C_0$ and $C_1$ denote the total half-edge configuration
weights of $B_0$ and $B_1$, respectively, then
$(3N/2-1)C_1=(3N/2)C_0$.  This identity has been verified only for
$N\le14$; the preceding proposition proves only its asymptotic consequence.
\end{conjecture}

This is a genuine reversible Markov-conductance consequence.  It is distinct
from GDPK's directed outer-conductance, whose volume is cardinality times the
constant outgoing slot count.

\begin{proposition}[Increasingly many vanishing Rule-A eigenvalues]
For every fixed $k$, the $k$th eigenvalue of the reversible normalized
Rule-A Laplacian $I-P_A$ satisfies
\[
    \lambda_k(I-P_A)\le \frac{C_k}{N}
\]
for all sufficiently large even $N$.
\end{proposition}

\begin{proof}
Choose $k-1$ distinct fixed cycle lengths.  Under the cubic configuration
model, the corresponding cycle-count variables converge jointly to
independent Poisson variables.  Hence the centered indicators that a carrier
contains a cycle of each chosen length have norms bounded away from zero and
an asymptotically nonsingular Gram matrix.

If a carrier contains a chosen $j$-cycle, a one-step move which destroys the
last such cycle must be performed on or within bounded distance of that cycle.
Only $O_j(1)$ of the $3N$ compressed slots can do this.  The Rule-A Dirichlet
energy of its centered indicator is therefore $O_j(N^{-1})$.  Cauchy--Schwarz
gives the same order for cross energies.  The min--max principle on the span of
these indicators and the constant function gives the result.
\end{proof}

This is the normalized multiplicity-chain analogue of Li's bounded-eigenvalue
construction.  It concerns the Markov operator actually sampled by local
Rule A, rather than either simple-support Laplacian.

The key distinction is:
\[
    \text{GoG bottleneck}
    \quad\neq\quad
    \text{bad field-averaging graph}.
\]
Grave de Peralta--Kolpakov's bridged set is a bottleneck in $\G_N$ because few
Whitehead moves leave the set of bridged graphs.  But the field relaxation at a
time $t$ is controlled by $\lambda_2(L_{G_t})$ or by the weighted analogue
coming from the edge-selection rule.  A graph with a microscopic bridge may
still have a spectral gap bounded below by a constant.  A balanced bridge is
bad for field averaging, because it gives a cut of size one separating two
linear-size pieces and hence $\lambda_2(L_G)=O(1/N)$ by the Rayleigh quotient.

More explicitly, if a bridge separates a set $S$ of $k$ vertices from its
complement, the test function which is constant on $S$ and on $S^c$ gives
\[
    \lambda_2(L_G)
    \le
    \frac{N}{k(N-k)}.
\]
This is $O(1/N)$ for balanced bridges and $O(1)$ for microscopic bridges.  Thus
bridged graphs are not uniformly bad for field relaxation; balanced sparse cuts
are.

The converse direction is governed by ordinary spectral graph theory.  Define
the edge isoperimetric constant of the current cubic graph by
\[
    h(G)=
    \min_{0<|S|\le N/2}
    \frac{|\partial_G S|}{|S|},
\]
where $\partial_G S$ is the set of edge objects crossing from $S$ to $S^c$.
For graphs of maximum degree $3$, Cheeger inequalities imply, up to universal
constants \cite{chung},
\[
    h(G)^2 \lesssim \lambda_2(L_G)\lesssim h(G).
\]
Thus a large sparse cut supplies a slowly contracting field direction, not a
statement that every current field contracts slowly.  The actual one-step
loss is governed by the Rayleigh quotient $f^TL_Gf/\V(f)$; a field nearly
orthogonal to the low mode may contract quickly.  Microscopic bridges can
coexist with $h(G)\asymp1$ when all cuts separating a linear-size set have
linear boundary, whereas balanced bridges force $h(G)=O(1/N)$.  This is the
precise interface between bridge geometry and the Lyapunov estimate.

The Rule-A stationary probability of a balanced cut of any fixed width
consequently tends to zero.  The bridged GoG bottleneck remains important for
the mixing of the base chain, but it does not bias its equilibrium carrier
marginal away from the configuration model.

\begin{conjecture}[Field exponent scale, after excluding killing words]
For a precisely specified Rule-A cocycle with no finite consensus word, or
along sizes for which such a word is impossible, the finite top exponent
satisfies
\[
    N\Lambda_N\to -\theta
\]
for some $\theta>0$.  The value may differ from the independent reset constant
because the evolving field and carrier are correlated.  Without the stated
qualification, the conjecture is false as a finite-valued assertion at
$N=2,4,8$, and possibly at further powers of two.
\end{conjecture}

\section{Results for the three graphs of moves}

The preceding field analysis uses the multiplicity-first Rule-A clock.  The
three base graphs also admit sharper, mutually incompatible conclusions.

\subsection{The GDPK multiplicity chain}

\begin{theorem}[Spectral interval, Perron root, and history entropy]
Let $M_N=(m_{GH})$ be the $3N$-slot adjacency matrix and
$P_A=M_N/(3N)$.  Then
\[
    \Spec(P_A)\subset[-1/2,1],
    \qquad \rho(M_N)=3N.
\]
In particular $M_N$ is real-diagonalizable, and from every source class the
number of concrete length-$t$ slot histories is exactly $(3N)^t$.
Furthermore the stationary entropy rate of the class-valued chain is
\[
 h_N=\log(3N)-\Delta_N,
 \qquad
 \Delta_N=\frac1{3N}\sum_G\pi_A(G)\sum_Hm_{GH}\log m_{GH}.
\]
If $C_1(G)=\sum_H\binom{m_{GH}}2$, then
\[
 0\le\Delta_N\le\frac{2}{3N}\E_{\pi_A}C_1(G).
\]
\end{theorem}

\begin{proof}
Detailed balance symmetrizes $P_A$.  Upstairs, freeze a central nonloop edge
and the four surrounding labelled stubs.  Their six two--two assignments form
the Johnson graph $J(4,2)$, whose normalized adjacency spectrum is
$\{1,0,-1/2\}$.  A loop gives the identity.  Averaging these local kernels and
restricting to class-invariant functions proves the spectral interval.  The
row sum of $M_N$ is exactly $3N$, proving the Perron and history claims.  The
entropy formula follows by expanding the entropy of the row probabilities
$m_{GH}/(3N)$; $\Delta_N$ is precisely the conditional entropy of the chosen
slot given source and target classes.  Finally,
$m\log m\le m(m-1)=2\binom m2$ for every positive integer $m$, which proves
the collision bound.
\end{proof}

\begin{theorem}[Exact loop-count Dirichlet form]
For $N\ge4$, let $L(G)$ be the number of carrier loops and
$\mu_N(k)=\pi_A\{L=k\}$.  Then for every function $f$ on the nonnegative
integers,
\[
 \mathcal E_A(f(L),f(L))
 =\frac{2}{3N}\sum_{k\ge1}k\mu_N(k)
       \bigl(f(k)-f(k-1)\bigr)^2.
\]
Consequently, for every fixed $r\ge1$,
\[
 \lambda_{r+1}(I-P_A)\le\frac{2r}{3N}+o(N^{-1}).
\]
With the convention
$\alpha_{\rm LS}=\inf\mathcal E(F,F)/\operatorname{Ent}_{\pi_A}(F^2)$,
also $\alpha_{\rm LS}\le2/(3N)+o(N^{-1})$.
\end{theorem}

\begin{proof}
A state with $k$ loops has exactly $2k$ compressed slots that destroy one
loop.  No move changes loop count by more than one when $N\ge4$.  Detailed
balance supplies the reverse aggregate flow and gives the displayed form.
Under $\pi_A$, $L$ tends to $\operatorname{Pois}(1)$.  Applying min--max to
the first $r$ Charlier polynomials, whose limiting immigration--death
Rayleigh quotients are $1,\ldots,r$, proves the eigenvalue bound.  The
indicator of $L=0$ proves the log-Sobolev upper bound.
\end{proof}

The projected process $L(G_t)$ need not itself be Markov.  The theorem is an
exact stationary Dirichlet-form identity and an upper bound, not convergence
of the low spectrum to a one-dimensional birth--death chain.

\subsection{The fixed-labelled simple graph}

Write $\Omega_N^{\rm lab}$ for connected adjacency multigraphs on the fixed
set $[N]$, $d(G)$ for Li's simple degree, and $\delta(G)=6N-d(G)$.  Three
natural walks have different stationary laws:
\[
 \begin{array}{c|c}
 \text{walk}&\text{stationary law}\\ \hline
 \text{simple neighbour}&d(G)/\sum_Hd(H)\\
 \text{uniform clock }I-L(\Gamma_{\rm lab})/(6N)&\text{uniform on }\Omega_N^{\rm lab}\\
 \text{four-slot multiplicity}&c(G)/\sum_Hc(H).
 \end{array}
\]
Denote the uniform law by $U_N$ and the four-slot stationary law by
$\pi_{\rm slot}$.

For connected $N>2$, put $L(G)=$ number of loops and $D(G)=$ number of
doubled vertex pairs.  Triple edges cannot occur, and exactly
\[
    c(G)=\frac{6^N}{2^{L(G)+D(G)}}.
\]
Under the slot law, $(L,D)$ tends to two independent
$\operatorname{Pois}(1)$ variables; exponential tilting shows that under the
uniform adjacency law it tends to two independent
$\operatorname{Pois}(2)$ variables.  Hence
\[
 \norm{\pi_{\rm slot}-U_N}_{\rm TV}
 \longrightarrow
 \norm{\operatorname{Pois}(1)^{\otimes2}
       -\operatorname{Pois}(2)^{\otimes2}}_{\rm TV}
 =0.4385731106\ldots.
\]
Li's good-edge estimate gives the stronger statement
$\E_{U_N}\delta=O(1)$, and therefore the simple-neighbour stationary law is
$O(N^{-1})$ in total variation from $U_N$.  Thus the two natural labelled
walks remain order-one apart even though their rows are typically
$O(N^{-1})$ apart.

More precisely, let $P_U=I-L(\Gamma_{\rm lab})/(6N)$, let $P_{\rm simp}$ be
the simple-neighbour kernel, and let $K(G,H)$ be the multiplicity among the
$6N$ labelled slots.  Put $P_{\rm slot}(G,H)=K(G,H)/(6N)$.  With
$k_0(G)=K(G,G)$,
\[
 \delta(G)=k_0(G)+\sum_{H\ne G:K(G,H)>0}(K(G,H)-1),
\]
and direct summation gives
\[
 \norm{P_{\rm simp}(G,\cdot)-P_U(G,\cdot)}_{\rm TV}
   =\frac{\delta(G)}{6N},\qquad
 \norm{P_{\rm slot}(G,\cdot)-P_U(G,\cdot)}_{\rm TV}
   \le\frac{\delta(G)}{6N}.
\]
These local estimates do not imply close stationary laws because the
relaxation scale is itself of order $N$.

\begin{theorem}[Quadratically many labelled low modes]
For
\[
 W_N=\left\{(a_{ij}):a_{ij}=a_{ji},\ a_{ii}=0,\
                 \sum_{j\ne i}a_{ij}=0\right\},
 \qquad \dim W_N=\frac{N(N-3)}2,
\]
put $F_a(G)=\sum_{i<j}a_{ij}m_{ij}(G)$.  There is an absolute $C$ such that,
for all sufficiently large even $N$,
\[
 \lambda_{N(N-3)/2+1}\bigl(L(\Gamma_{\rm lab}(N))\bigr)\le C.
\]
The same number of nonconstant modes have eigenvalue $O(N^{-1})$ for the
uniform-clock and normalized simple-neighbour operators.
\end{theorem}

\begin{proof}
The $N$ row-sum equations have full rank, which gives the displayed dimension,
and $W_N$ is the irreducible $[N-2,2]$ representation of $S_N$.  Both the
uniform covariance of $F_a$ and its simple-edge Dirichlet form are
$S_N$-invariant, so Schur's lemma makes each a scalar form on $W_N$.

Let $x(G)=(m_{ij}(G))_{i<j}$.  The trace of the covariance form on $W_N$ is
$\E_{U_N}\|\operatorname{Proj}_{W_N}x(G)\|^2$.  Since
$\sum_{i<j}m_{ij}^2=3N/2+O(L+D)$ in expectation, while the trivial and
standard projections have bounded expected squared norm, division by
$\dim W_N$ gives
\[
 \E_{U_N}F_a^2=\left(\frac3N+O(N^{-2})\right)\norm{a}^2.
\]
Each raw move changes only four pair coordinates.  Summing the resulting
squared increments over slots and taking the trace on $W_N$ gives simple-edge
Dirichlet energy $O(N^{-1})\norm{a}^2$; invariance again makes the same bound
uniform over $a\in W_N$.  Its quotient by the covariance is $O(1)$, while the
uniform-clock normalization contributes an additional factor $(6N)^{-1}$.

For the normalized simple walk the denominator is the invariant form
$\E_{U_N}[d(G)F_a(G)^2]$.  Its trace is $\Theta(N^2)$ because
$d=6N-\delta$, $\E_{U_N}\delta=O(1)$, and
$\|\operatorname{Proj}_{W_N}x\|^2=O(N)$ deterministically.  Hence this form is
$\Theta(1)\norm{a}^2$ and the normalized quotient is $O(N^{-1})$.  Adding the
constant function and applying min--max proves all three assertions.
\end{proof}

\begin{theorem}[Fixed-radius labelled balls]
Let $G\sim U_N$.  For every fixed $R$,
\[
 |B_R^{\rm lab}(G)|=\frac{(6N)^R}{R!}+O_{\mathbb P,R}(N^{R-1}),
 \qquad
 \E_{U_N}|B_R^{\rm lab}(G)|
 =\frac{(6N)^R}{R!}+O_R(N^{R-1}).
\]
The error estimate is not asserted uniformly for growing $R$.
\end{theorem}

\begin{proof}
Call a raw move $R$-good when the radius-$(2R+2)$ carrier neighbourhood of its
central edge has no loop or parallel edge and is a tree rooted at that edge.  A
word is separated if all of these neighbourhoods
are vertex-disjoint.  Standard fixed-subgraph counting in the cubic
configuration model \cite{janson-random-graphs}, followed by the exact
loop/double-edge exponential tilt
to $U_N$, preserves these fixed-radius moment bounds and shows that the
expected number of non-good raw slots is
$O_R(1)$.  Once fewer than $R$ separated supports have been chosen, only
$O_R(1)$ of the $6N$ next raw slots meet one of them.  Hence the expected number
of ordered separated good words is
\[
 (6N)^R+O_R(N^{R-1}),
\]
and the same estimate holds in probability.

The rewrites in a separated word commute.  In the symmetric difference of the
source and endpoint edge sets, the affected neighbourhoods are distinct
components; within each, the unchanged central edge is the unique diagonal.
Thus the source--endpoint pair recovers the $R$ selected moves and their
resolutions.  Exactly the $R!$ orders of those commuting moves have the same
endpoint.  Words with repeated or intersecting supports, or with a non-good
slot, have expected count $O_R(N^{R-1})$, because one free $O(N)$ choice is
lost.  Words of length less than $R$ have the same lower order.  Dividing the
good ordered count by $R!$ proves the expectation estimate, and Markov's
inequality applied to the exceptional-word count gives the in-probability
form.
\end{proof}

Moreover
\[
 |\Omega_N^{\rm lab}|=(e^2+o(1))\frac{(3N-1)!!}{6^N},
\]
so maximum degree $6N$ implies
$\diam\Gamma_{\rm lab}(N)\ge(3/2-o(1))N$, and the same lower bound holds for
the distance between two uniform labelled states with high probability.

\begin{proposition}[Exact orbit quotient]
For an unlabelled class $\alpha$, put
$s_\alpha=N!/|\Aut_v(\alpha)|$, and let $b_{\alpha\beta}$ be the number of
labelled simple neighbours of one representative of $\alpha$ lying in orbit
$\beta$.  Then
\[
 s_\alpha b_{\alpha\beta}=s_\beta b_{\beta\alpha}.
\]
The orbit partition is strongly lumpable for the labelled simple-neighbour
walk, with transition $b_{\alpha\beta}/d_\alpha$ and stationary mass
proportional to $s_\alpha d_\alpha$.  The uniform-clock walk lumps with
off-diagonal and diagonal transitions
\[
 \overline P_U(\alpha,\beta)=\frac{b_{\alpha\beta}}{6N}
 \quad(\beta\ne\alpha),\qquad
 \overline P_U(\alpha,\alpha)
 =\frac{6N-d_\alpha+b_{\alpha\alpha}}{6N},
\]
and stationary mass proportional to $s_\alpha$.  The unlabelled simple support records only whether
$b_{\alpha\beta}>0$ and is not, in general, either Markov quotient.
\end{proposition}

\begin{proof}
The symmetric group acts transitively on every orbit and by automorphisms of
the labelled simple graph, so $b_{\alpha\beta}$ is representative-independent.
Double-counting labelled edges between the two orbits gives reciprocity.
Summing its rows gives the stated lumped transition and stationary laws.
Discarding the values of $b_{\alpha\beta}$ retains reachability but not these
probabilities.
\end{proof}

\subsection{The unlabelled simple support}

Reciprocity makes the support undirected.  With
$q(G)=\#\{H\ne G:m_{GH}>0\}$, simple-neighbour Rule B has
\[
 P_B(G,H)=\frac1{q(G)},\qquad
 \pi_B(G)=\frac{q(G)}{\sum_Kq(K)}.
\]
Condition Rule A on a nonself move and put
\[
 M_+(G)=\sum_{H\ne G}m_{GH},\qquad
 C_1^+(G)=\sum_{H\ne G}\binom{m_{GH}}2.
\]
Then the discrepancy from a uniform support neighbour is exactly
\[
 \chi^2(p_A^+\Vert p_B)
 =\frac{q(G)(M_+(G)+2C_1^+(G))}{M_+(G)^2}-1.
\]
Thus the collision statistic controls both support degree and row-law
discrepancy.

If $C_1(G)=\sum_H\binom{m_{GH}}2$ includes the self target, Cauchy--Schwarz
also gives the exact deterministic lower bound
\[
 q(G)\ge
 \frac{(3N)^2}{3N+2C_1(G)}-1.
\]
Indeed, at most $q(G)+1$ targets, including the self target, have positive
multiplicity, while
\[
 \sum_Hm_{GH}^2=3N+2C_1(G).
\]

For a class set $X$, let $b_X$ be its simple-support boundary and
$V_X=\sum_{G\in X}q(G)$.  Its centered indicator has exact Rule-B quotient
\[
 R_B(X)=b_X\left(\frac1{V_X}+\frac1{V_{X^c}}\right).
\]
\begin{proposition}[Conditional Rule-B low modes]
Suppose a fixed short-cycle event $X$ and its complement both have cardinality
$\Omega(|\C_N|)$, and all but $o(|\C_N|)$ unlabelled classes have support
degree at least $d_N\to\infty$.  Then $\lambda_2(I-P_B)\to0$.  More generally,
for a fixed $k$, suppose there is a family of $k-1$ short-cycle events whose
centered indicators have an asymptotically nonsingular class-uniform Gram
matrix.  If every atom of that family has conditional mean degree $\Omega(N)$,
then
\[
 \lambda_k(I-P_B)=O_k(N^{-1}).
\]
These conclusions are conditional: the stated degree hypotheses remain open.
\end{proposition}

\begin{proof}
For the positive-density short-cycle feature $X$, Li's local argument gives
$b_X\le C_X|X|$.  The first hypothesis makes both $V_X$ and $V_{X^c}$
at least a constant times $d_N|\C_N|$, so the exact quotient is
$O_X(d_N^{-1})$; min--max with the constant function gives the first claim.
For finitely many features, their centered indicators have an asymptotically
nonsingular Gram matrix by hypothesis, while their
diagonal and cross Dirichlet energies are $O_k(N^{-1})$ after normalization.
Min--max on their span gives the fixed-index assertion.
\end{proof}

In particular, two-sided conditional estimates
$\E[C_1(G)\mid G\in X]=O(N)$ and
$\E[C_1(G)\mid G\in X^c]=O(N)$ imply the required linear mean-degree bounds
for that feature by the collision inequality and Jensen's inequality.

\begin{proposition}[Active-edge field kernel]
Call a carrier edge active when at least one resolution has a nonself target.
Assume Rule B gives positive probability to every support neighbour and every
concrete lift of that transition.  Then
\[
 \ker L_G^B
 =\{f:f\text{ is constant on every component of the active-edge carrier}\}.
\]
The weighted gap vanishes exactly when that carrier is disconnected, with
mean-zero kernel dimension one less than its component count.  The obstruction
is instantaneous: subsequent rewrites can change the active components.
\end{proposition}

\begin{proof}
Under the full-support lift rule, an edge has positive averaging weight exactly
when it is active.  Since
\[
 f^TL_G^Bf=\sum_e p_G^B(e)(f_i-f_j)^2,
\]
the quadratic form vanishes precisely when $f_i=f_j$ along every active edge,
which is equivalent to constancy on each active component.  Intersecting this
kernel with $\one^\perp$ removes one dimension.
\end{proof}

Finally, multiplicities and self-slots do not change reachability.  Hence the
simple support has exactly the GDPK support metric and inherits its
$\Theta(N\log N)$ diameter order \cite{rafi-tao}, even though its path counts
and random-walk times differ.

\section{Complete finite enumeration}

The finite statements in this section are computational evidence, not
asymptotic theorems.  Cubic multigraphs were encoded by coloured incidence
graphs and canonically labelled with nauty/Traces \cite{mckay-nauty}.
Enumeration is complete through $N=14$: the class counts agree with
A005967, and every source has exactly $3N$ compressed resolution slots.
\[
\begin{array}{c|r|r}
N&\text{unlabelled classes}&\text{compressed slots}\\ \hline
2&2&12\\
4&5&60\\
6&17&306\\
8&71&1\,704\\
10&388&11\,640\\
12&2\,592&93\,312\\
14&21\,096&886\,032
\end{array}
\]

\subsection{Carrier and base-chain spectra}

The minimum carrier gap over all classes is compared below with the universal
path bound $a_N=2(1-\cos(\pi/N))$.
\[
\begin{array}{c|r|r|r}
N&\min_G\lambda_2(L_G)&a_N&\min_G\lambda_2(L_G)/a_N\\ \hline
2&2.000000&2.000000&1.000\\
4&0.763932&0.585786&1.304\\
6&0.354249&0.267949&1.322\\
8&0.202067&0.152241&1.327\\
10&0.130145&0.097887&1.330\\
12&0.090687&0.068148&1.331\\
14&0.066764&0.050144&1.331
\end{array}
\]
At every enumerated $N\ge4$, the minimizer is the alternating
single/double-edge path with endpoint loops used in the proof of the
$O(N^{-2})$ extremal upper bound.

The first six eigenvalues of the Rule-A kernel, Rule-B simple-neighbour
kernel, and unnormalized support Laplacian are:
\[
\begin{array}{c|c|rrrrrr}
N&\text{operator}&1&2&3&4&5&6\\ \hline
8&I-P_A&.000&.046&.097&.124&.161&.194\\
 &I-P_B&.000&.074&.163&.221&.275&.314\\
 &L_{\rm supp}&.000&.307&.617&.935&1.048&1.153\\
10&I-P_A&.000&.031&.063&.082&.087&.107\\
 &I-P_B&.000&.046&.099&.135&.163&.187\\
 &L_{\rm supp}&.000&.307&.504&.757&.870&.999\\
12&I-P_A&.000&.023&.045&.060&.065&.077\\
 &I-P_B&.000&.030&.065&.092&.105&.117\\
 &L_{\rm supp}&.000&.335&.682&.879&1.057&1.135\\
14&I-P_A&.000&.018&.034&.046&.050&.056\\
 &I-P_B&.000&.021&.045&.066&.070&.080\\
 &L_{\rm supp}&.000&.332&.706&.788&.862&1.116
\end{array}
\]
These values are consistent with the proved upper bounds on fixed-index
normalized eigenvalues.  They do not prove a matching lower bound.

\subsection{Bridge cut and Rule-B field obstruction}

The bridge-cut conductances are
\[
\begin{array}{c|r|r|r}
N&\text{support boundary edges}&\Phi_A(B_N)&\Phi_B(B_N)\\ \hline
8&15&.0909&.2113\\
10&82&.0714&.1673\\
12&537&.0588&.1221\\
14&4\,425&.0500&.0922
\end{array}
\]
The Rule-A column agrees exactly with $1/(3N/2-1)$ for every even
$N\le14$.  This finite identity is a conjecture beyond the enumerated
range; only its asymptotic consequence
$\Phi_A(B_N)=2/(3N)(1+o(1))$ is proved.

Pure nonself Rule B has a disconnected active-edge carrier on the following
classes:
\[
\begin{array}{c|r|r|r|r}
N&\text{zero-gap classes}&\text{class fraction}
 &\pi_A\text{-mass}&\pi_B\text{-mass}\\ \hline
8&17&.239&.139&.215\\
10&89&.229&.101&.195\\
12&548&.211&.0757&.174\\
14&4\,102&.194&.0616&.1565
\end{array}
\]
Every such class in the enumeration is bridged.  The active-edge theorem
makes the zero instantaneous gaps exact; the table does not imply a zero
trajectory Lyapunov exponent.

\subsection{Fixed-labelled graph}

The labelled simple graph was materialized completely for $N\le6$.
\[
\begin{array}{c|r|r|r|r|r}
N&\text{states}&\text{simple edges}&\E_Ud&
  \min d/\max d&\lambda_2(L_{\rm lab})\\ \hline
2&2&1&1.0000&1/1&2.000000\\
4&35&99&5.6571&5/6&2.052917\\
6&3\,550&23\,400&13.1831&9/24&1.646846
\end{array}
\]
For larger $N$, exact orbit weighting gives
\[
\begin{array}{c|r|r}
N&\E_Ud&6N\\ \hline
8&22.232&48\\
10&32.105&60\\
12&42.501&72\\
14&53.266&84
\end{array}
\]
The sizes shown are not yet in the asymptotic $6N-O(1)$ regime, but the
ratios increase monotonically over this range.  Exact rational arithmetic
also verifies the $K_4$ and $Q_3$ consensus products, and direct evaluation
of the resolution-independent energy identity on sampled transitions agrees
to machine precision.

\paragraph{Reproducibility.}
The script \path{gog_generate.py} constructs the canonical class and
transition tables.  The reported multiplicity/support spectra and bridge
statistics are produced by \path{gog_edge_model_experiments.py};
\path{gog_support_feature_analysis.py} evaluates short-cycle and Rule-B
diagnostics, and \path{gog_ball_growth.py} samples support balls.
Labelled-degree and automorphism checks are implemented in
\path{gog_degree_diagnostics.py}.  Exact commands, conventions, and
numerical limitations are recorded in \path{docs/computation.rst}.

\section{Lean Verification Boundary}

A self-contained Lean~4 project in \texttt{lean/} kernel-checks the exact
finite algebra that is actually formalized.  It uses exact rational arithmetic
and the standard library only; the project source contains no \texttt{sorry},
\texttt{admit}, or project-declared \texttt{axiom}.  Representative
\texttt{\#print axioms} checks report only Lean's foundational principles
\texttt{propext}, \texttt{Classical.choice}, and \texttt{Quot.sound}.  The
checked declarations include:
\begin{itemize}[leftmargin=*]
    \item the coordinate definition of pair averaging, symmetry, affine and
    centering covariance, preservation of the selected-pair sum, idempotence,
    and the exact squared-norm/variance drop $(x_i-x_j)^2/2$ with an arbitrary
    unchanged remainder;
    \item commutation of averages on disjoint pairs;
    \item identity, composition, and inverse laws for pullback by an explicit
    coordinate permutation, including exact conjugacy between relabelling and
    pair averaging;
    \item preservation of total mass, commutation and idempotence of the block
    schedules, and exact consensus for the two transverse $K_4$ matchings and
    the three coordinate matchings of $Q_3$;
    \item a common-denominator proof that no finite pair-average schedule sends
    the three-state datum $(0,0,1)$ to exact consensus;
    \item concrete $2\times2$ identities $A^2=A$, $A\one=\one$, and agreement
    between matrix and coordinate pair averaging.
\end{itemize}
From the \texttt{lean/} directory, \texttt{lake build} checks the project.
The theorem map is in \path{lean/README.md}.

This is deliberately not advertised as a formalization of the whole paper.
Lean does not presently check the Whitehead graph or its quotient, construct
settlements, prove their coherence, or address stabilizer holonomy.  Nor does
it check the configuration counts, the representation-theoretic argument,
random-graph inputs, asymptotics, or Lyapunov and mixing theorems.  The dyadic
result is a concrete three-state obstruction rather than a classification of
finite consensus protocols.  The formal layer certifies only the finite
averaging identities on which several of those arguments rely.

\section{Open Problems}

\begin{question}
After fixing a settlement cocycle and excluding sizes with finite killing
words, is the top Lyapunov exponent asymptotic to $-\theta/N$?  If yes, is
$\theta$ determined only by typical cubic expansion, or does it remember
carrier--field alignment and projection noncommutativity?
\end{question}

\begin{question}
Can the square-defect curvature be related to low GoG Laplacian modes?  The
plausible mechanism is that local motifs responsible for bottleneck eigenvectors
also create high curvature density for incident Whitehead moves, but this is a
testable statement rather than a theorem.
\end{question}

\begin{question}
Does the finite bridge identity
$(3N/2-1)C_1=(3N/2)C_0$ hold for all even $N$, and can one prove a matching
$c/N$ Poincar\'e lower bound for the labelled local Johnson-update lift?
\end{question}

\begin{question}
Can accidental isomorphisms among local rewrites be controlled strongly enough
to prove growing conditional degree in the unlabelled simple support?  A
linear bound would turn the conditional Rule-B low-mode statements above into
unconditional $O(N^{-1})$ spectral bounds.
\end{question}

\begin{question}
Is there a useful Hodge-like theory for the covariant difference $D$ on the
field bundle over $\G_N$?  Such a theory must first specify edge-to-two-cell
differentials for a noninvertible semigroup cocycle and treat automorphism
groups as finite gauge stabilizers.
\end{question}

\begin{thebibliography}{12}

\bibitem{gdp-kolpakov}
Laura Grave de Peralta and Alexander Kolpakov,
\emph{Expansion properties of Whitehead moves on cubic graphs},
Comptes Rendus Mathématique \textbf{362} (2024), 1825--1836.
\url{https://doi.org/10.5802/crmath.691}

\bibitem{li}
Michael Li,
\emph{Increasingly Many Bounded Eigenvalues of the Graph of Whitehead Moves},
arXiv:2405.14592.
\url{https://arxiv.org/abs/2405.14592}

\bibitem{oeis-a005967}
OEIS Foundation Inc.,
\emph{A005967: Number of isomorphism classes of connected 3-regular multigraphs of order $2n$, loops allowed},
The On-Line Encyclopedia of Integer Sequences.
\url{https://oeis.org/A005967}

\bibitem{mckay-nauty}
Brendan D. McKay and Adolfo Piperno,
\emph{Practical graph isomorphism, II},
Journal of Symbolic Computation \textbf{60} (2014), 94--112.
\url{https://doi.org/10.1016/j.jsc.2013.09.003}

\bibitem{chung}
Fan R. K. Chung,
\emph{Spectral Graph Theory},
CBMS Regional Conference Series in Mathematics, vol. 92,
American Mathematical Society, 1997.
\url{https://bookstore.ams.org/cbms-92/}

\bibitem{janson-random-graphs}
Svante Janson, Tomasz {\L}uczak, and Andrzej Ruci\'nski,
\emph{Random Graphs},
Wiley Series in Discrete Mathematics and Optimization,
John Wiley \& Sons, 2000.

\bibitem{friedman}
Joel Friedman,
\emph{A Proof of Alon's Second Eigenvalue Conjecture and Related Problems},
Memoirs of the American Mathematical Society \textbf{195} (2008), no. 910.
\url{https://bookstore.ams.org/MEMO/195/910}

\bibitem{bordenave}
Charles Bordenave,
\emph{A New Proof of Friedman's Second Eigenvalue Theorem and Its Extension to Random Lifts},
Annales Scientifiques de l'\'Ecole Normale Sup\'erieure (4) \textbf{53}
(2020), no.~6, 1393--1439.
\url{https://doi.org/10.24033/asens.2450}

\bibitem{rafi-tao}
Kasra Rafi and Jing Tao,
\emph{The Diameter of the Thick Part of Moduli Space and Simultaneous
Whitehead Moves},
Duke Mathematical Journal \textbf{162} (2013), no.~10, 1833--1876.
\url{https://doi.org/10.1215/00127094-2323128}

\bibitem{levin-peres-wilmer}
David A. Levin and Yuval Peres,
\emph{Markov Chains and Mixing Times},
2nd ed., with contributions by Elizabeth L. Wilmer,
American Mathematical Society, 2017.
\url{https://pages.uoregon.edu/dlevin/MARKOV/mcmt2e.pdf}

\bibitem{williams}
David Williams,
\emph{Probability with Martingales},
Cambridge University Press, 1991.
\url{https://doi.org/10.1017/CBO9780511813658}

\bibitem{durrett}
Rick Durrett,
\emph{Probability: Theory and Examples},
5th ed., Cambridge University Press, 2019.
\url{https://doi.org/10.1017/9781108591034}

\end{thebibliography}

\end{document}
